\documentclass{article}
\usepackage{../Self_Style}

\title{UChicago REU Notes: Monads and Monadicity}
\author{Zih-Yu Hsieh}
\date{\today}

\begin{document}
\maketitle
\tableofcontents

\section{Review: Adjoint Functors}

This part is converging to the preparation of the theory of monads :) (as professor May had mentioned it quite often as an essential tool to understand operad, so I decided to do related review). The main reference is Riehl's \emph{Category Theory in Context}.

\begin{defn}
    Let $\mathcal{C},\mathcal{D}$ be two categories. Two opposite functors $F:\mathcal{C}\rightarrow \mathcal{D}$, $\mathcal{C}\leftarrow \mathcal{D}:G$ are \emph{adjoint pairs} - in particular, $F$ is a \emph{left adjoint of $G$},and $G$ is a \emph{right adjoint of $F$} - if there is a natural isomorphism
    \[\Hom_\mathcal{D}(F(X),Y)\cong \Hom_\mathcal{C}(X,G(Y))\]
    For any $X\in\mathcal{C}$ and $Y\in\mathcal{D}$.

    Notation wise, it's denoted as $(F,G)$ or $F\dashv G$.
\end{defn}

More formally, it's a natural isomorphism $(\_)^\sharp:\Hom_\mathcal{D}(F(\_),\_)\xrightarrow{\sim}\Hom_\mathcal{C}(\_,G(\_))$ of bifunctors $\mathcal{C}^\Op\times\mathcal{D}\rightarrow\Sets$ (here, we also denote $(\_)^\flat$ as the inverse transformation). 

This means, given morphism $f:X\rightarrow X'$ in $\mathcal{C}$, and morphism $g:Y\rightarrow Y'$ in $\mathcal{D}$, we have diagrams as follow:
\[\xymatrix{
    \Hom_\mathcal{D}(F(X'),Y) \ar[d]_-{\_\circ Ff}\ar[r]^-{\sim} & \Hom_\mathcal{C}(X',G(Y)) \ar[d]^-{\_\circ f}\\
    \Hom_\mathcal{D}(F(X),Y)\ar[r]^-{\sim} & \Hom_\mathcal{C}(X,G(Y))
},\quad \quad \xymatrix{
    \Hom_\mathcal{D}(F(X),Y) \ar[d]_-{g\circ \_} \ar[r]^-{\sim} & \Hom_\mathcal{C}(X,G(Y)) \ar[d]^-{Gg\circ \_}\\
    \Hom_\mathcal{D}(F(X),Y')\ar[r]^-{\sim} & \Hom_\mathcal{C}(F(X),Y')
}\]
The first diagram insists that any morphism $h:F(X')\rightarrow Y$ in $\mathcal{D}$ satisfies $(h\circ Ff)^\sharp = h^\sharp \circ f$ in $\mathcal{C}$. The second diagram insists that any morhpism $\ell:F(X)\rightarrow Y$ satisfies $(g\circ \ell)^\sharp = Gg\circ \ell^\sharp$. Or, the following diagram holds in $\mathcal{C}$:
\[\xymatrix{
    X' \ar[d]_-{f} \ar[dr]^-{(h\circ f)^\sharp}\\
    X \ar[r]_-{h^\sharp} & G(Y)
},\quad\quad \xymatrix{
    X \ar[r]^-{\ell^\sharp} \ar[dr]_-{(g\circ \ell)^\sharp} & G(Y) \ar[d]^-{Gg}\\
    & G(Y')
}\]
If use the inverse transformation $(\_)^\flat$, the above can also be phrased as: Any morphism $\tilde{h}:X'\rightarrow G(Y)$ satisfies $(\tilde{h}\circ f)^\flat = (\tilde{h})^\flat \circ Ff$, and any morhpism $\tilde{\ell}:X\rightarrow G(Y)$ satisfies $(Gg\circ \tilde{\ell})^\flat = g\circ (\tilde{\ell})^\flat$ in $\mathcal{D}$.

\hfil

Adjoint functor also connects the commutativity of certain diagrams between the two categories:

\begin{lem}
Fix arbitrary $f:X\rightarrow X'$ in $\mathcal{C}$, $g:Y\rightarrow Y'$ in $\mathcal{D}$, and two morphsisms $h:F(X)\rightarrow Y$, $h':F(X')\rightarrow Y'$ in $\mathcal{D}$. The commutativity of following two diagrams are equivalent:
\[\xymatrix{
    F(X) \ar[d]_-{Ff} \ar[r]^-{h} & Y \ar[d]^-{g}\\
    F(X')\ar[r]_-{h'} & Y'
} \text{ in }\mathcal{D},\quad \quad \xymatrix{
    X \ar[d]_-{f} \ar[r]^-{h^\sharp} & G(Y) \ar[d]^-{Gg}\\
    X'\ar[r]_-{h'^\sharp} & G(Y')
} \text{ in }\mathcal{C}\]
\end{lem}
\begin{proof}
    As a remark, $(\_)^\sharp$ is a bijection on each homset, so two morphisms $F(X)\rightarrow Y'$ are equal iff after applying $(\_)^\sharp$ they're equal as morphisms $X\rightarrow G(Y')$. 
    
    Using the notation from above, we see the following equivalence:
    \[g\circ h=h'\circ Ff \iff (g\circ h)^\sharp = (h'\circ Ff)^\sharp \iff Gg\circ h^\sharp = h'^\sharp \circ f \]
    Where, the second equivalence is the property from above. The left and right most equalities represent the two commutative diagrams respectively, hence showing the equivalence of commutativity.
\end{proof}

\hfil

Given the adjoint properties, we have the following two natural isomorphisms, given any $X\in\mathcal{C}$ and $Y\in\mathcal{D}$:
\[(\_)^\sharp:\Hom_\mathcal{D}(F(X),F(X))\xrightarrow{\sim} \Hom_\mathcal{C}(X,GF(X)),\quad\quad \Hom_\mathcal{D}(FG(Y),Y)\xleftarrow{\sim} \Hom_\mathcal{C}(G(Y),G(Y)):(\_)^\flat\]
Hence, we can define a collection of maps $\eta_X := (\id_{F(X)})^\sharp:X\rightarrow GF(X)$, and $\epsilon_Y := (\id_{G(Y)})^\flat:FG(Y)\rightarrow Y$.

\begin{prop}
The two collections form natural transformations $\eta:\Id_\mathcal{C}\rightarrow GF$, and $\epsilon:FG\rightarrow \Id_\mathcal{D}$.
\end{prop}
\begin{proof}
    Here, fix arbitrary $f:X\rightarrow X'$ in $\mathcal{C}$, and $g:Y\rightarrow Y'$ in $\mathcal{D}$. Using \textbf{Lemma 1.5}, the following two diagrams have equivalent commutativity:
    \[\xymatrix{
        F(X) \ar[d]_-{Ff} \ar[r]^-{\id_{F(X)}} & F(X) \ar[d]^-{Ff}\\
        F(X)\ar[r]_-{\id_{F(X')}} & F(X')
    } \text{ in }\mathcal{D},\quad \quad \xymatrix{
        X \ar[d]_-{f} \ar[rr]^-{(\id_{F(X)})^\sharp=\eta_X} && GF(X) \ar[d]^-{GFf}\\
        X'\ar[rr]_-{(\id_{F(X')})^\sharp=\eta_{X'}} && GF(X')
    } \text{ in }\mathcal{C}\]
    Where, the left side obviously commutes, so the right side commutes, showing naturality of $\eta$.

    Dually, since $(\epsilon_Y)^\sharp = \id_{G(Y)}$ by the inverse property, we also have the equivalence of the following two diagrams:
    \[\xymatrix{
        FG(Y) \ar[d]_-{FGg} \ar[r]^-{\epsilon_Y} & Y \ar[d]^-{g}\\
        FG(Y)\ar[r]_-{\epsilon_{Y'}} & Y'
    } \text{ in }\mathcal{D},\quad \quad \xymatrix{
        G(Y) \ar[d]_-{f} \ar[rr]^-{(\epsilon_Y)^\sharp = \id_{G(Y)}} && G(Y) \ar[d]^-{Gg}\\
        G(Y')\ar[rr]_-{(\epsilon_{Y'})^\sharp=\id_{G(Y')}} && G(Y')
    } \text{ in }\mathcal{C}\]
    Where, the right side obviously commutes, so the left side commutes, showing the naturality of $\epsilon$.
\end{proof}

These natural transformations are named as follow:
\begin{defn}
    Given the adjoint pairs $F\dashv G$, the natural transformation $\eta:\Id_\mathcal{C}\rightarrow GF$ is called the \emph{unit transformation}, while $\epsilon:FG\rightarrow \Id_\mathcal{D}$ is called the \emph{counit transformation}.
\end{defn}

\hfil

They satisfy what's called the \emph{Triangle Identity} as follow:
\begin{prop}
    The following two diagrams of natural transformation holds:
    \[\xymatrix{
        F \ar[r]^-{F\eta} \ar[dr]_-{\id_F} & FGF \ar[d]^-{\epsilon_F}\\
        & F
    } \text{ in }\mathcal{D}^\mathcal{C},\quad \quad \xymatrix{
        G \ar[r]^-{\eta_G} \ar[dr]_-{\id_G} & GFG \ar[d]^-{G\epsilon}\\
        & F
    } \text{ in }\mathcal{C}^\mathcal{D}\]
\end{prop}
Here, $\mathcal{D}^\mathcal{C}$ denotes the \emph{functor categories} of all functors $\mathcal{C}\rightarrow \mathcal{D}$, with morphisms being natural transformations between functors.

\begin{proof}
    Given arbitrary $X\in\mathcal{C}$, the goal is to show the horizontal morphisms of the following diagram are identities:
    \[\xymatrix{
        F(X) \ar[r]^-{F\eta_X} & FGF(X) \ar[r]^-{\epsilon_{F(X)}} & F(X)
    }\]
    Using out notation and the naturality, we see the following equality of morphisms after applying $(\_)^\sharp$:
    \[(\epsilon_{F(X)}\circ F\eta_X)^\sharp = (\epsilon_{F(X)})^\sharp \circ \eta_X = \id_{GF(X)}\circ \eta_X = \eta_X:GF(X)\rightarrow X\]
    But, recall that $\eta_X:= (\id_{F(X)})^\sharp$, using the bijectivity of $(\_)^\sharp$, this enforces $\epsilon_{F(X)}\circ F\eta_X = \id_{F(X)}$, showing $\epsilon_F\circ F\eta=\id_F$ as natural transformations.

    \hfil

    Dually, given arbitrary $Y\in \mathcal{D}$, the goal is again to show the following morphism is identity:
    \[\xymatrix{
        G(Y) \ar[r]^-{\eta_{G(Y)}} & GFG(Y) \ar[r]^-{G\epsilon_{Y}} & G(Y)
    }\]
    Similarly, after applying $(\_)^\flat$ we get:
    \[(G\epsilon_Y\circ \eta_{G(Y)})^\flat = \epsilon_Y\circ (\eta_{G(Y)})^\flat = \epsilon_Y\circ \id_{FG(Y)}= \epsilon_Y\]
    The second equality relies on the fact $(\id_{F(X)})^\sharp = \eta_X$. Which, since $(\_)^\flat$ is again a bijection, and $\epsilon_Y := (\id_G(Y))^\flat$, we see $G\epsilon_Y\circ \eta_{G(Y)}=\id_{G(Y)}$, showing $G\epsilon\circ \eta_{G} = \id_G$ as natural transformations.
\end{proof}

This will be an important result when classifying monads. 

\hfil

We also want to say the converse of Triangle Identity holds, that it generates an adjoint pair:
\begin{lem}
    Suppose $F:\mathcal{C}\rightarrow \mathcal{D}$ and $\mathcal{C}\leftarrow \mathcal{D}:G$ are two opposite functors, together with natural transformations $\eta:\Id_\mathcal{C}\rightarrow GF$ and $\epsilon:FG\rightarrow \Id_\mathcal{D}$ such that Triangle Identity holds:
    \[\xymatrix{
        F \ar[r]^-{F\eta} \ar[dr]_-{\id_F} & FGF \ar[d]^-{\epsilon_F}\\
        & F
    } \text{ in }\mathcal{D}^\mathcal{C},\quad \quad \xymatrix{
        G \ar[r]^-{\eta_G} \ar[dr]_-{\id_G} & GFG \ar[d]^-{G\epsilon}\\
        & F
    } \text{ in }\mathcal{C}^\mathcal{D}\]
    then $F\dashv G$ is an adjoint pair.
\end{lem}
\begin{proof}
    Given any $X\in \mathcal{C}$ (associated to $\eta_X:X\rightarrow GF(X)$) and $Y\in\mathcal{D}$ (associated to $\epsilon_Y:FG(Y)\rightarrow Y$), for any morphism $h:F(X)\rightarrow Y$, we have a corresponding morphism $Gh:GF(X)\rightarrow G(Y)$. This defines the following map:
    \[(\_)^\sharp:\Hom_\mathcal{D}(F(X),Y)\rightarrow \Hom_\mathcal{C}(X,G(Y)),\quad h^\sharp := Gh\circ \eta_X\]
    Similarly, for any morphism $\ell:X\rightarrow G(Y)$, we have a corresponding morphism $F\ell:F(X)\rightarrow FG(Y)$. This again defines the following map:
    \[\Hom_\mathcal{D}(F(X),Y)\leftarrow \Hom_\mathcal{C}(X,G(Y)):(\_)^\flat,\quad \ell^\flat := \epsilon_Y\circ F\ell\]
    These are evidently functorial. To show they're mutual inverses, we expand the following:
    \[(h^\sharp)^\flat = \epsilon_Y\circ Fh^\sharp = \epsilon_Y\circ F(Gh\circ \eta_X) = \epsilon_Y \circ FGh\circ F\eta_X\]
    Expanding the right hand side in terms of diagram (includes the triangle identity and naturality of $\epsilon$), we get:
    \[\xymatrix{
        F(X) \ar@(u,u)[rr]^-{\id_{F(X)}} \ar[r]^-{F\eta_X} & FGF(X) \ar[d]_-{FGh} \ar[r]^-{\epsilon_{F(X)}} & F(X)\ar[d]^-{h}\\
        & FG(Y) \ar[r]_-{\epsilon_{Y}} & Y
    }\]
    So, we see $(h^\sharp)^\flat = \epsilon_Y\circ FGh\circ F\eta_X = h$ by the above commutative diagram.

    Expand the dual diagram with the other triangle identity and naturality of $\eta$, we'll get that $(\ell^\flat)^\sharp = \ell$ instead, showing the two are mutual inverses.
\end{proof}

\newpage

\section{Example: Free-Forgetful Adjunction Monad}

\newpage

\section{Monads and Relations to Adjoint Functors}

Let's first observe the case when we have an adjoint pair $F\dashv G$ from above, inducing an endofunctor $GF:\mathcal{C}\rightarrow\mathcal{C}$, with unit transformation $\eta:\Id_\mathcal{C}\rightarrow GF$, together with \emph{multiplication transformation} $\mu:=G\epsilon_F:GFGF\rightarrow GF$ (as we have $\epsilon:FG\rightarrow\Id_\mathcal{D}$). Notice they provide the following identities:
\begin{itemize}
    \item \emph{Associativity of Multiplication}:
    \[\xymatrix{
        (GF)^3 \ar[d]_-{\mu_{GF}} \ar[r]^-{GF\mu} & (GF)^2\ar[d]^-{\mu}\\
        (GF)^2 \ar[r]_-{\mu} & GF
    }\]
    Fix arbitrary $X\in\mathcal{C}$, expand the definition, we get the following diagram:
    \[\xymatrix{
        GFG(FGF(X)) \ar[d]_-{G\epsilon_{FGF(X)}} \ar[r]^-{GFG\epsilon_{F(X)}} & GFG(F(X)) \ar[d]^-{G\epsilon_{F(X)}}\\
        G(FGF(X)) \ar[r]_-{G\epsilon_{F(X)}} & G(F(X))
    }\]
    Disregard the left most $G$, the above diagram commutes because $\epsilon:FG\rightarrow\Id_\mathcal{D}$ is a natural transformation.

    \item \emph{Unit Property}:
    \[\xymatrix{
        GF \ar[dr]_-{\id_{GF}} \ar[r]^-{\eta_{GF}} & (GF)^2 \ar[d]^-{\mu} & GF \ar[l]_-{GF\eta} \ar[dl]^-{\id_{GF}}\\
        & GF
    }\]
    Using the definition $\mu:= G\epsilon_F$, we have the following two equalities using \emph{Triangle Identities of Adjunction}:
    \[G\epsilon_F \circ \eta_{GF} = (G\psi_F\circ \eta_G)_F = (\id_G)_F = \id_{GF}\]
    \[G\epsilon_F\circ GF\eta = G(\epsilon_F \circ F\eta) = G(\id_F) = \id_{GF} \]
    Thse two identities reflect the commutativity of the above diagram.
\end{itemize}

This reflects the structure of \emph{Monad}:
\begin{defn}
    Given a category $\mathcal{C}$, a \emph{Monad} on $\mathcal{C}$ consists of the following:
    \begin{itemize}
        \item An endofunctor $T:\mathcal{C}\rightarrow\mathcal{C}$.
        \item A \emph{unit transformation} $\eta:\Id_\mathcal{C}\rightarrow T$.
        \item A \emph{multiplication transformation} $\mu:T^2\rightarrow T$.
    \end{itemize}
    There are extra requirements of satisfying the following two diagrams as endofunctors on $\mathcal{C}$:
    \[\xymatrix{
        T^3 \ar[d]_-{\mu_{T}} \ar[r]^-{T\mu} & T^2\ar[d]^-{\mu}\\
        T^2 \ar[r]_-{\mu} & T
    },\quad\quad \xymatrix{
        T \ar[dr]_-{\id_{T}} \ar[r]^-{\eta_{T}} & T^2 \ar[d]^-{\mu} & T \ar[l]_-{T\eta} \ar[dl]^-{\id_{T}}\\
        & T
    }\]
    Which are analogs to \emph{associativity} and \emph{unit} of monoids.
\end{defn}

Which, starting portion shows that any adjoint pair $(F\dashv G)$ generates a monad on category $\mathcal{C}$ (the source category of the left adjoint $F$).

\begin{comment}

\begin{exmp}
    As a more relevant example, let $\mathcal{T}^*$ denotes the category of based spaces, $\mathcal{T}$ denotes the category of topological spaces, there is a natural forgetful functor $F:\mathcal{T}^*\rightarrow \mathcal{T}$ that forgets the based spaces.

    It has a left adjoint $(\_)_*:\mathcal{T}\rightarrow\mathcal{T}^*$, where $(X)_* := X\sqcup \{*\}$ by adding a point as basepoint, and any map $f:X\rightarrow Y$ naturally extends to a based map $\tilde{f}:X_*\rightarrow Y_*$ where $\tilde{f}|_X = f$, and $\tilde{f}(*)=*$.

    For any space $X\in\mathcal{T}$ and based space $Y\in\mathcal{T}^*$ (with $y_0\in Y$ as basepoint), the natural isomorphism $\Hom_{\mathcal{T}^*}(X_*, Y)\cong \Hom_\mathcal{T}(X,Y)$ is given by $g\mapsto g|_X$, and inverse is $f\mapsto \tilde{f}$.

    \hfil

    In this case, the unit map $\eta_X = \id_{X_*}|_X$, which becomes the canonical inclusion $X\hookrightarrow X_*$; the counit map $\epsilon_Y = \tilde{\id_Y}$, which is the map $Y_*\rightarrow Y$ that collapse the new basepoint $*$ to the original basepoint $y_0$. So, the multiplication map $\eta_X = F\epsilon_{X_*}:(X_*)_*\rightarrow X_*$ by collapsing the two added basepoints (in different iteration) to the unique basepoint (with one iteration).

    Since collapsing of three added points to one doesn't care about the collapse order, this explicitly forms a monad (or using the adjunction property this directly follows :)
\end{exmp}

\end{comment}

\hfil

Another question to ask: Is every monad of a category induced by adjoint pairs of functors? The answer is yes, but one needs to introduce some relevant categories constructed out of the monads.

\begin{defn}
    Let $(T,\eta, \mu)$ be a monad on category $\mathcal{C}$. Its \emph{Eilenberg-Moore category} for $T$ (or the \emph{category of $T$-algebras}), is a category $\mathcal{C}^T$ consists of:
    \begin{itemize}
        \item Objects of pairs $(A,a)$, where $A\in\mathcal{C}$ is an object, and $a:T(A)\rightarrow A$ is a morphism, so the following diagram commutes:
        \[\xymatrix{
            A \ar[dr]_-{\id_A} \ar[r]^-{\eta_A} & T(A) \ar[d]^-{a}\\
            & A
        },\quad\quad \xymatrix{
            T^2(A) \ar[r]^-{\mu_A} \ar[d]_-{Ta} & T(A) \ar[d]^-{a}\\
            T(A) \ar[r]_-{a} & A
        }\]
        \item A morphism $f:(A,a)\rightarrow (B,b)$ is a morhpism $f:A\rightarrow B$, such that the following diagram commutes:
        \[\xymatrix{
            T(A) \ar[d]_-{a} \ar[r]^-{Tf} & T(B)\ar[d]^-{b}\\
            A \ar[r]_-{f} & B
        }\]
    \end{itemize}
\end{defn}

Intuitively $T(A)$ consists of data about ``pairs of actions'', while $a$ encodes how this action appears on $A$. The square provides the ``associativity'' of action (for instance, $(ab)\cdot v = a\cdot (b\cdot v)$ as action on modules). I do need to agree with professory May, that this system is more like a module instead of algebra in the traditional sense.

\hfil

For this category, there is an obvious ``forgetful functor'' $G^T:\mathcal{C}^T\rightarrow \mathcal{C}$ by $G^T (A,a) = A$ and $G^T f = f$ for any $f:(A,a)\rightarrow (B,b)$. 

Similarly, since $\mathcal{C}$ is given a monad, there is also an obvious ``algebra functor'' $F^T:\mathcal{C}\rightarrow \mathcal{C}^T$, by $F^T(A) := (T(A),\mu_A)$, and $F^Tf = Tf$ for any $f:A\rightarrow B$. This is because $\mu_A$ satisfies the following two identities, by the definition of monad:
\[\xymatrix{
    T(A) \ar[dr]_-{\id_{T(A)}} \ar[r]^-{\eta_{T(A)}} & T^2(A) \ar[d]^-{\mu_{T(A)}}\\
    & T(A)
},\quad\quad \xymatrix{
    T^2(T(A)) \ar[r]^-{\mu_{T(A)}} \ar[d]_-{T\mu_A} & T(T(A)) \ar[d]^-{\mu_A}\\
    T(T(A)) \ar[r]_-{\mu_A} & T(A)
}\]
Also, since $\mu:T^2\rightarrow T$ is a natural transformation, the following diagram commutes:
\[\xymatrix{
    T(T(A)) \ar[d]_-{\mu_A} \ar[r]^-{Tf} & T(T(B))\ar[d]^-{\mu_B}\\
    T(A) \ar[r]_-{f} & T(B)
}\]
showing $Tf:(T(A),\mu_A)\rightarrow (T(B),\mu_B)$ is a morphism.

Here comes an important result regarding this category:
\begin{thm}
    The two functors are adjoint pairs $(F^T\dashv G^T)$, and the induced monad is $(T,\eta,\mu)$.
\end{thm}%reference


This demonstrates each monad is induced by some adjoint pairs.

\begin{comment}

\begin{exmp}
    For the basepoint monad demonstrated in \textbf{Example 1.11}, the Eilenberg-Moore category $\mathcal{T}^{(\_)_*}$ consists of pairs $(X,a)$ where $a:X_*\rightarrow X$ is a map that satisfies the following:
    \[\xymatrix{
        X \ar[dr]_-{\id_X} \ar[r]^-{\eta_X} & X_* \ar[d]^-{a}\\
        & X
    },\quad\quad \xymatrix{
        (X_*)_* \ar[r]^-{Ta} \ar[d]_-{\mu_X} & X_* \ar[d]^-{a}\\
        X_* \ar[r]_-{a} & X
    }\]
    The first diagram enforces $a|_X = \id_X$ (as $\eta_X:X\rightarrow X_*$ is the canonical inclusion $X\hookrightarrow X\sqcup\{*\}$). Given this data, the second diagram automatically commutes (by checking where the added points are sent to), so we see $a$ is purely determined by $a(*)\in X$, in other words, it's specified by $(X,x_0)$ for $x_0\in X$.

    Then, given a morphism $f:(X,a)\rightarrow (Y,b)$, the following diagram commutes:
    \[\xymatrix{
        X_* \ar[d]_-{a} \ar[r]^-{\tilde{f}} & Y_* \ar[d]^-{b}\\
        X \ar[r]_-{f} & Y
    }\]
    Which, $f(a(*))= b(\tilde(f)(*))=b(*)$ is enforced, showing $f$ must map the chosen point in $X$ to the chosen point in $Y$.

    Hence, the functor $\mathcal{T}^{(\_)_*}\rightarrow \mathcal{T}^*$ by $(X,a)\mapsto (X,a(x))$, and $f\mapsto f$ is in fact an equivalence of category. So, the Eilenberg-Moore category of the basepoint functor is the category of based space.
\end{exmp}
\end{comment}

\hfil

Afterward, there's more question to ask: Given a monad on $\mathcal{C}$, how many ways are there to create adjoint pairs which induce the provided monad? Can we find such universal adjoint pair? 
\begin{thm}
    In the \emph{adjoint category} $\text{Adj}_T$ (categories of all category with adjoint functors between $\mathcal{C}$, inducing monad $T$; the morphisms are functors that commutes with the left/right adjoints for each category), the Eilenberg-Moore category $(\mathcal{C}^T,F^T\dashv G^T)$ is final.
\end{thm}

\newpage

\section{Monad Algebra and Monadicity}

Fix a monad $(T,\eta,\mu)$ on category $\mathcal{C}$, we're mainly concerned with the Eilenber-Moore category $\mathcal{C}^T$ (or the category of $T$-algebras). In particular, we wish to study the free algebras, namely the pair $(T(A),\mu_A)\in \mathcal{C}^T$, and see under what situation this covers all the possible $T$-algebras.

The first goal is to show in $\mathcal{C}^T$, any $T$-algebra $(A,a)$ satisfies the following coequalizer diagram:
\[\xymatrix{
    (T^2(A),\mu_{T(A)}) \ar@/_{-.5pc}/[r]^-{Ta} \ar@/_{.5pc}/[r]_-{\mu_A} & (T(A),\mu_A) \ar[r]^-{a} & A
}\]

For this, let's do some definition:
\begin{defn}
    A \emph{split coequalizer diagram} consists of the following morphisms:
    \[\xymatrix{
        X \ar@/_{-.5pc}/[r]^-{f} \ar@/_{.5pc}/[r]_-{g} & Y \ar@(d,d)[l]^-{t} \ar[r]^-{h} & Z \ar@(d,d)[l]^-{s}
    }\]
    with the requirement $h\circ f=h\circ g$, $h\circ s=\id_Z$, $g\circ t = \id_Y$, and $f\circ t=s\circ h$ as endomorphism on $Y$.
\end{defn}

\begin{lem}
    Given a split coequalizer diagram as above, then $h:Y\rightarrow Z$ serves as a categorical coequalizer of $f,g:X\rightarrow Y$. Also, it's an \emph{absolute colimit}, meaning all functor preserves this coequalizer.
\end{lem}
\begin{proof}
    It is evident why all functor preserves this, simply because functor preserves identity, hence split coequalizer diagram is preserved under any functor. Hence, it suffices to show a split coequalizer diagram is a coequalizer.

    Let $h':Y\rightarrow Z'$ satisfies $h'\circ f=h'\circ g$, consider the morphism $h'\circ s:Z\rightarrow Y\rightarrow Z'$. Notice its composition with $h$ provides the following:
    \[(h'\circ s)\circ h = h'\circ (f\circ t) = (h'\circ g)\circ t = h'\circ \id_{Y} = h'\]
    Hence, the following factorization holds:
    \[\xymatrix{
        X \ar@/_{-.5pc}/[r]^-{f} \ar@/_{.5pc}/[r]_-{g} & Y\ar[r]^-{h} \ar[dr]_-{h'} & Z \ar[d]^-{h'\circ s}\\
        && Z'
    }\]
    Then, the uniqueness of $h'\circ s$ is due to the fact $h$ is an epimorphism: Recall that $h\circ s=\id_Z$, with $\id_Z$ being epimorphism, so is $h$. Then, if $\psi:Z\rightarrow Z'$ also satisfies $(h'\circ s)\circ h=h' = \psi\circ h$, right cancellative property guarantees $\psi=h'\circ s$. This shows $h:Y\rightarrow Z$ is a coequalizer of $f,g$.
\end{proof}

For instance, given $(A,a)\in \mathcal{C}^T$ a $T$-algebra, the following diagram is a split coequalizerin $\mathcal{C}$ (\textbf{Remark:} NOT in $\mathcal{C}^T$):
\[\xymatrix{
    T^2(A) \ar@/_{-.5pc}/[r]^-{Ta} \ar@/_{.5pc}/[r]_-{\mu_A} & T(A) \ar@(d,d)[l]^-{\eta_{T(A)}} \ar[r]^-{a} & A \ar@(d,d)[l]^-{\eta_A}
}\]
Since $a\circ Ta = a\circ \mu_A$ using associativity diagram while $a\circ \eta_A = \id_A$ by unit diagram of monad algebra, $\mu_A\circ \eta_{T(A)}=\id_{T(A)}$ by triangle identity of monad, and finally $Ta \circ \eta_{T(A)}=\eta_A\circ a$ due to the fact $\eta:\Id_\mathcal{C}\rightarrow T$ is a natural transformation.

\hfil

More general result need more definition:
\begin{defn}
    Given a functor $U:\mathcal{D}\rightarrow\mathcal{C}$.
    \begin{itemize}
        \item a \emph{$U$-split coequalizer} of $f,g:X\rightarrow Y$ in $\mathcal{D}$, is an object $Z\in\mathcal{C}$, together with morphism $h:U(Y)\rightarrow Z$, $s:Z\rightarrow U(Y)$, and $t:U(Y)\rightarrow U(X)$, such that the following diagram is a split coequalizer in $\mathcal{C}$:
        \[\xymatrix{
            U(X) \ar@/_{-.5pc}/[r]^-{Uf} \ar@/_{.5pc}/[r]_-{Ug} & U(Y) \ar@(d,d)[l]^-{t} \ar[r]^-{h} & Z \ar@(d,d)[l]^-{s}
        }\]
        \item $U$ \emph{creates coequalizers of $U$-split pairs}, if any $U$-split coequalizer admits a coequalizer $h':Y\rightarrow Z'$ in $\mathcal{D}$, such that its image $Uh':U(Y)\rightarrow U(Z')$ is isomorphic to $h:U(Y)\rightarrow Z$ in $\mathcal{C}$.
        \item $U$ \emph{strictly creates coequalizers of $U$-split pairs}, if it creates coequalizer of $U$-split pairs, while such lift $h'$ in $\mathcal{D}$ is unique (up to isomorphism in $\mathcal{D}$).
    \end{itemize}
\end{defn}

With these definitions, let's consider the property possessed by forgetful functor $G^T:\mathcal{C}^T\rightarrow\mathcal{C}$ for Eilenberg-Moore category:
\begin{prop}
    $G^T$ strictly creates $G^T$-split coequalizers.
\end{prop}

\hfil

Finally, given adjoint pair $F:\mathcal{C}\leftrightarrows\mathcal{D}:G$ (where $(F\dashv G)$) that induces monad $T$ on $\mathcal{C}$. Recall, there exists a canonical functor $K:\mathcal{D}\rightarrow \mathcal{C}^T$, such that the adjoint pairs commute:
\[\xymatrix{
    \mathcal{D} \ar@/_{-.5pc}/[dr]^-{G} \ar@{-->}[rr]^-{K} && \mathcal{C}^T \ar@/_{-.5pc}/[dl]^-{G^T}\\
    & \mathcal{C} \ar@/_{-.5pc}/[ul]^-{F} \ar@/_{-.5pc}/[ur]^-{F^T}
}\]

Here comes the famous \emph{Beck's Monadicity Theorem}, classifying when $\mathcal{D}$ is isomorphic to $\mathcal{C}^T$:
\begin{thm}
    $K$ is an isomorphism of categories $\iff$ $G$ strictly creates $G$-split coequalizers.
\end{thm}

\newpage

\section{Monad associated to an Operad}

\newpage

\section{Appendix 1: Categorical Proofs}

\subsection{Proof of Theorem 3.3}
\begin{proof}
    Given arbitrary $A\in\mathcal{C}$ and $(B,b)\in\mathcal{C}^T$, we shall show a natural isomorphism between the two:
    \[\Phi:\Hom_{\mathcal{C}^T}((T(A),\mu_A), (B,b))\rightleftarrows \Hom_\mathcal{C}(A,B):\Psi\]
    Given a morphism $f:(T(A),\mu_A)\rightarrow (B,b)$, recall the existence of unit morphism $\eta_A:A\rightarrow T(A)$, so $\Phi(f):= f\circ \eta_A:A\rightarrow T(A)\rightarrow B$ defines a map. 
    
    Conversely, given any morhpism $g:A\rightarrow B$, it naturally promotes to $Tg:(T(A),\mu_A)\rightarrow (T(B),\mu_B)$; and, since $(B,b)\in\mathcal{C}^T$, postcomposing with the morphism $b:(T(B),\mu_B)\rightarrow (B,b)$ provides $b\circ Tg:(T(A),\mu_A)\rightarrow (T(B),\mu_B)\rightarrow  (B,b)$. Define $\Psi(g):= b\circ Tg$.

    \begin{rem}
        The reason why $b:T(B)\rightarrow B$ is also a morphism in $\mathcal{C}^T$, is because of the commutative square as follow by the fact $(B,b)\in\mathcal{C}^T$:
        \[\xymatrix{
            T(T(B)) \ar[r]^-{\mu_B} \ar[d]_-{Tb} & T(B) \ar[d]^-{b}\\
            T(B) \ar[r]_-{b} & B
        }\ \implies \ \xymatrix{
            T(T(B)) \ar[r]^-{Tb} \ar[d]_-{\mu_B} & T(B) \ar[d]^-{b}\\
            T(B) \ar[r]_-{b} & B
        }\]
        The right hand side is the transpose of the left hand side, also the defining diagram of $b:(T(B),\mu_B)\rightarrow (B,b)$ as a morphism in $\mathcal{C}^T$.
    \end{rem}

    The two maps $\Phi,\Psi$ are for sure functorial, hence suffices to check they're mutual inverses. Let's compute them respectively:
    \begin{itemize}
        \item \[\Psi(\Phi(f)) = \Psi(f\circ \eta_A) = b\circ T(f\circ \eta_A)= b\circ Tf\circ T\eta_A\]
        To reduce this, consider the following diagram, using triangle identities of monad and the assumptions about morphisms:
        \[\xymatrix{
            T(A) \ar[dr]_-{\id_{T(A)}} \ar[r]^-{T\eta_A}& T^2(A) \ar[d]_-{\mu_A} \ar[r]^-{Tf} & T(B) \ar[d]^-{b}\\
            & T(A) \ar[r]_-{f} & B
        }\]
        So, we see $b\circ Tf \circ T\eta_A = f\circ \id_{T(A)}=f:T(A)\rightarrow B$ as morphism in $\mathcal{C}$ (and hence in $\mathcal{C}^T$), showing $\Psi(\Phi(f))=f$.
        \item \[\Phi(\Psi(g)) = \Phi(b\circ Tg)= b\circ Tg \circ \eta_A\]
        To reduce this, again consider similar definitions (and the fact $\eta:\Id_\mathcal{C}\rightarrow T$ is a natural transformation):
        \[\xymatrix{
            A \ar[r]^-{g} \ar[d]_-{\eta_A} & B \ar[d]_-{\eta_B} \ar[dr]^-{\id_B}\\
            T(A) \ar[r]_-{Tg} & T(B) \ar[r]_-{b} & B
        }\]
        So, we get $b\circ Tg\circ \eta_A = \id_B\circ g=g$ as morphism in $\mathcal{C}$, hence $\Phi(\Psi(g))=g$.
    \end{itemize}
    This shows $(F^T\dashv G^T)$ in fact forms an adjoint pair.

    \hfil

    To show they generate monad $(T,\eta,\mu)$ on $\mathcal{C}$, notice the following action on each object $A\in\mathcal{C}$ / morphism $f:A\rightarrow B$:
    \[G^T F^T(A) = G^T(T(A),\mu_A) = T(A),\quad G^T F^T f = G^T(T f) = Tf\]
    Hence, $G^T F^T = T$ as endofunctor on $\mathcal{C}$.

    Similarly, the unit transformation $\eta':\Id_\mathcal{C}\rightarrow G^T F^T$ as $\eta'_A := \Phi(\id_{(T(A),\mu_A)}) = \id_{T(A)}\circ \eta_A = \eta_A$, showing $\eta'=\eta$.

    Finally, we have counit transformation $\epsilon':F^TG^T \rightarrow \Id_{\mathcal{C}^T}$ as $\epsilon'_{(B,b)}:= \Psi(\id_B) = b\circ T(\id_B) = b\circ \id_{T(B)}=b$. So, the induced multiplication transformation $\mu' := G^T \epsilon'_{F^T}:(G^T F^T)^2\rightarrow G^T F^T$ satisfies $\mu'_A = G^T \epsilon'_{(T(A),\mu_A)} = G^T(\mu_A) = \mu_A$, showing $\mu'=\mu$.

    Hence, $(F^T\dashv G^T)$ induces monad $(T,\eta,\mu)$.
\end{proof}

\hfil

\subsection{Proof of Theorem 3.4}
\begin{proof}
    Given a category $\mathcal{D}$, together with adjoint pairs $(F\dashv G)$ (where $F:\mathcal{C}\leftrightarrows\mathcal{D}:G$) that generates monad $(T,\eta,\mu)$. We ought to show the existence / uniqueness of a functor $L:\mathcal{D}\rightarrow\mathcal{C}^T$, such that the following diagram commutes:
    \[\xymatrix{
        \mathcal{D} \ar@/_{.5pc}/[dr]_{G} \ar@{-->}[rr]^-{\exists ! L} && \mathcal{C}^T \ar@/_{.5pc}/[dl]_{G^T}\\
        & \mathcal{C} \ar@/_{.5pc}/[ul]_{F} \ar@/_{.5pc}/[ur]_{F^T}
    }\]
    For the $G$-diagram to commute, $L$ realizes a $T$-algebra structure on $G(Y)$ for each $Y\in\mathcal{D}$. Which, the counit $\epsilon:FG\rightarrow \Id_{\mathcal{D}}$ induces a natural transformation $\theta=G\epsilon:(GF)G\rightarrow G$, so each $Y\in\mathcal{D}$ is associated with a morphism $\theta_Y:T(G(Y))\rightarrow G(Y)$.

    Notice $(G(Y),\theta_Y)$ forms a $T$-algebra due to one of the Triangle Identity, and the ``reduced'' associativity diagram for adjunction (by swapping the last $F(X)$ with $Y$):
    \[\xymatrix{
        G \ar[r]^-{\eta_G} \ar[dr]_-{\id_G} & GFG \ar[d]^-{G\epsilon}\\
        & F
    },\quad\quad \xymatrix{
        GFGFG(Y) \ar[d]_-{G\epsilon_{FG(Y)}} \ar[r]^-{GFG\epsilon_{Y}} & GFG(Y) \ar[d]^-{G\epsilon_{Y}}\\
        GFG(Y) \ar[r]_-{G\epsilon_{F(X)}} & G(Y)
    }\]
    Also, by naturality of $\epsilon$, given any morphism $f:Y\rightarrow Y'$ in $\mathcal{D}$, $Gf:G(Y)\rightarrow G(Y')$ in $\mathcal{C}$ forms a morphism of $T$-algebra. Hence, define $L:\mathcal{D}\rightarrow \mathcal{C}^T$ by $L(Y):= (G(Y), G\epsilon_Y)$, and $Lf:= Gf:G(Y)\rightarrow G(Y')$ is a functor.

    This functor satisfies $G(Y) = G^T L(Y)$, and $Gf = G^T Lf$, because $L(Y) = (G(Y),G\epsilon_Y)$ (also $Lf = Gf$), while forgetful functor $G^T(G(Y),G\epsilon_Y) = G(Y)$ (and morphism $L(Gf)= Gf$). 
    
    Also, given $X\in\mathcal{C}$ and morphism $g:X\rightarrow X'$ in $\mathcal{C}$, we have $L(F(X)) = (GF(X), G\epsilon_{F(X)}) = (T(X),\mu_X) = F^T(X)$ (by checking the definition of multiplication transformation, and the fact $GF = T$); also, $L(Ff) = G(Ff):GF(X)\rightarrow GF(X')$ by definition, which equals to $F^T f:T(X)\rightarrow T(X')$, this shows $LF = F^T$ as functor. This shows the existence of $L$.

    \hfil

    The uniqueness of $L$ on the morphism is clear, as $L(Y)\in\mathcal{C}^T$ must have the underlying object being $G(Y)\in\mathcal{C}$ (for $G^TL = G$ to be true), so $f:Y\rightarrow Y'$ must be mapped to $Gf:G(Y)\rightarrow G(Y')$. Theuniqueness of $L$ on the object is by the monadic requirement, where the $T$-algebra structure on $G(Y)$ is restricted by the monad requirement.
\end{proof}

\hfil

\subsection{Proof of Proposition 4.4}
\begin{proof}
    Given two $T$-algebras $(A,\alpha),(B,\beta)$ with two morphisms $f,g:(A,\alpha)\rightrightarrows(B,\beta)$ in $\mathcal{C}^T$ that admits a $G^T$-splitting:
    \[\xymatrix{
            A \ar@/_{-.5pc}/[r]^-{f} \ar@/_{.5pc}/[r]_-{g} & B \ar@(d,d)[l]^-{t} \ar[r]^-{h} & C \ar@(d,d)[l]^-{s}
        }\]
    Recall that $h\circ f=h\circ g$, together with $h\circ s=\id_C$, $g\circ t=\id_B$, and $f\circ t=s\circ h$ as endomorphism of $B$. Here, the goal is to show, $C$ has a unique $T$-algebra structure $\gamma:T(C)\rightarrow C$, such that $h$ is the unique morphism forming $(C,\gamma)$ as a qoequalizer of $f,g$ in $\mathcal{C}^T$.
    
    Using the fact that $C$ (with morphism $h$) is a coequalizer in $\mathcal{C}$, together with the fact $T$ preserves split coequalizer (in fact, all functor preserves split coequalizer), we see $Th:T(B)\rightarrow T(C)$ is the coequalizer of $Tf,Tg:T(A)\rightarrow T(B)$. Also, since $f,g$ are algebra morphisms, we get the following diagram:
    \[\xymatrix{
        T(A) \ar[d]_-{\alpha} \ar@/_{-.5pc}/[r]^-{Tf} \ar@/_{.5pc}/[r]_-{Tg} & T(B) \ar[d]_-{\beta} \ar[r]^-{Th} & T(C) \ar@{-->}[d]^-{\exists ! \gamma} \\
        A \ar@/_{-.5pc}/[r]^-{f} \ar@/_{.5pc}/[r]_-{g} & B \ar[r]^-{h} & C 
        }\]
        Existence and uniqueness of $\gamma$ is based on the fact $h\circ \beta$ coequalizes $Tf,Tg$, and $T(C)$ with $Th$ forms a coequalizer.

        \hfil

        So, the goal is to prove $(C,\gamma)$ forms a $T$-algebra and is a coequalizer of $f,g$ in $\mathcal{C}^T$, as existence/uniqueness (up to isomorphism) of the lift follows from the universality of coequalizer.

        To check $(C,\gamma)$ is a $T$-algebra, we use the fact $(B,\beta)$ is a $T$-algebra, together with the fact $h$ is an epimorphism (as a coequalizer), with the ``algebra morphism structure'' above (Remark: Here $h$ is not an algebra morphism yet, it's just a morphism that satisfies the same diagram for algebra morphism at this point). Which, the following diagrams commute:
        \[\xymatrix{
            & B \ar[rr]^-{h} \ar[dl]_-{\eta_B} \ar[dd]_-{\id_B} && C \ar[dl]_-{\eta_C} \ar[dd]^-{\id_C}\\
            T(B) \ar[rr]^-{Th} \ar[dr]_-{\beta} && T(C) \ar[dr]_-{\gamma}\\
            & B \ar[rr]_-{h} && C
        },\quad\quad \xymatrix{
            & T^2(B) \ar[rrr]^-{T^2h} \ar[dl]_-{T\beta} \ar[ddr]_-{\mu_B} &&& T^2(C) \ar[dl]_-{T\gamma} \ar[ddr]^-{\mu_C}\\
            T(B) \ar[rrr]^-{Th} \ar[ddr]_-{\beta} &&& T(C) \ar[ddr]_{\gamma}\\
            && T(B) \ar[dl]_-{\beta} \ar[rrr]^-{Th} &&& T(C) \ar[dl]^-{\gamma}\\
            & B \ar[rrr]_-{h} &&& C
        }\]
        Where, for both diagram, the right side are the axioms for $T$-algebra, while the other sides of the diagrams commute by checking the naturality / algebra. 

        Using the fact $T$ preserves split coequalizer (and $h$ is a split coequalizer), we deduce $T^2h$ is also a split coequalizer. So, by chasing the diagram we get:
        \[\id_C\circ h=h\circ \id_B = h\circ \beta\circ \eta_B = \gamma\circ Th\circ \eta_B = \gamma\circ \eta_C\circ h\]
        \[\gamma\circ \mu_C\circ T^2h = \gamma\circ Th\circ \mu_B=h\circ \beta\circ \mu_B=h\circ \beta\circ T\beta=\gamma\circ Th\circ T\beta = \gamma\circ T\gamma \circ T^2h\]
        Using the fact $h,T^2h$ are coequalizers (hence epimorphisms), we get $\gamma\circ \eta_C = \id_C$, and $\gamma\circ \mu_C = \gamma\circ T\gamma$, hence $(C,\gamma)$ is a $T$-algebra, further implying $h:B\rightarrow C$ is also a $T$-algebra morphism $(B,\beta)\rightarrow (C,\gamma)$.

        \hfil

        Finally, $h$ is again a coequalizer in $\mathcal{C}^T$, because any coequalizing diagram in $\mathcal{C}^T$ forgets to be a coequalizing diagram in $\mathcal{C}$, hence using the fact $h$ is a coequalizer in $\mathcal{C}$, the factorization follows (and the remaining part is to check it's a morphism of $T$-algebra, which follows from the fact that $T$ preserves split coequalizer).
\end{proof}

\newpage

\section{Appendix 2: Kleisli Category}

Given a monad $T$ on $\mathcal{C}$, there is another special category with adjoint pairs that induces the monad $T$:
\begin{defn}
    Given a monad $(T,\eta,\mu)$ on category $\mathcal{C}$. The \emph{Kleisli category} $\mathcal{C}_T$ consists of:
    \begin{itemize}
        \item Objects are carried over from $C$.
        \item Morphism $f:A\rightsquigarrow B$ in $\mathcal{C}_T$ is a morphism $f:A\rightarrow T(B)$ in $\mathcal{C}$. Under this definition, $\eta_A:A\rightsquigarrow A$ serve as the identity morphism of $A\in\mathcal{C}_T$, while given $f:A\rightsquigarrow B$ and $g:B\rightsquigarrow C$, the following morphism in $\mathcal{C}$ serves as the composition of the two:
        \[g\circ_* f:= A \xrightarrow{f} T(B) \xrightarrow{Tg} T(T(C)) \xrightarrow{\mu_C}T(C)\]
    \end{itemize}
\end{defn}
\begin{proof}
    First, let's show associativity: Given $f:A\rightsquigarrow B$, $g:B\rightsquigarrow C$, and $h:C\rightsquigarrow D$, the two ways of composition provide the following in $\mathcal{C}$:
    \[h\circ_* (g\circ_* f) = A\xrightarrow{f}T(B)\xrightarrow{Tg}T^2(C)\xrightarrow{\mu_C}T(C) \xrightarrow{Th}T^2(D)\xrightarrow{\mu_D}T(D)\]
    \[(h\circ_* g)\circ_* f = A\xrightarrow{f}T(B) \xrightarrow{Tg} T^2(C) \xrightarrow{T^2h}T^3(D) \xrightarrow{T\mu_D}T^2(D) \xrightarrow{\mu_D} T(D)\]
    To show the two procedures agree, it suffices to show the middle right portion $\mu_D\circ Th\circ \mu_C = \mu_D \circ T\mu_D \circ T^2h$. However, this follows by the monad property:
    \[\xymatrix{
        T^2(C) \ar[r]^-{T^2 h}\ar[d]_-{\mu_C} & T^3(D) \ar[d]^-{\mu_{T(D)}} \ar[r]^-{T\mu_{D}} & T^2(D) \ar[d]^-{\mu_D}\\
        T(C) \ar[r]_-{Th} & T^2(D) \ar[r]_-{\mu_D} & T(D)
    }\]
    the left square commutes by naturality of $\mu$, the right square commutes by associativity of $\mu$.

    \hfil

    To show $\eta_A$ serves as identity, we compute the following given $f:A\rightsquigarrow B$ and $g:Z\rightsquigarrow A$:
    \begin{itemize}
        \item \[f\circ_* \eta_A = A\xrightarrow{\eta_A}T(A) \xrightarrow{Tf} T^2(B) \xrightarrow{\mu_B}T(B)\]
        This can be simplified using triangle identity, and the fact $\eta:\Id_\mathcal{C}\rightarrow T$ is a natural transformation:
        \[\xymatrix{
            A \ar[d]_-{\eta_A} \ar[r]^-{f} & T(B) \ar[d]_-{\eta_{T(B)}} \ar[dr]^-{\id_{T(B)}}\\
            T(A) \ar[r]_-{Tf} & T^2(B) \ar[r]_-{\mu_{B}} & T(B)
        }\]
        So, we deduce $f\circ_* \eta_A =\mu_B\circ Tf\circ \eta_A = \id_{T(B)}\circ f = f$.
        \item \[\eta_A\circ_* g = Z\xrightarrow{g}T(A) \xrightarrow{T\eta_A} T^2(A) \xrightarrow{\mu_A}T(A)\]
        Notice that $\mu_A\circ T\eta_A = \id_{T(A)}$ in $\mathcal{C}$ (by triangle identity),so we see $\mu_A\circ_* T\eta_A =\mu_A\circ T\eta_A \circ g = \id_{T(A)}\circ g= g$.
    \end{itemize}
\end{proof}

\hfil

For this category, there is a functor $G_T:\mathcal{C}_T\rightarrow \mathcal{C}$, by $G_T (A) = T(A)$, and $G_T f:= \mu_B\circ Tf:T(A)\rightarrow T^2(B)\rightarrow T(B)$ for arbitrary $A\in\mathcal{C}_T$ and morphism $f:A\rightsquigarrow B$. This is functorial, as $G_T \eta_A = \mu_A\circ T\eta_A = \id_{T(A)}$ by applying triangle identity, and given $f:A\rightsquigarrow B$, $g:B\rightsquigarrow C$, we have the following diagram commutes:
\[\xymatrix{
    T(A) \ar[r]^-{Tf} & T^2(B) \ar[d]_-{\mu_B} \ar[r]^-{T^2 g} & T^3(C) \ar[d]_-{\mu_{T(C)}} \ar[r]^-{T\mu_C} & T^2(C) \ar[d]^-{\mu_C}\\
    & T(B) \ar[r]_-{Tg} & T^2(C) \ar[r]_-{\mu_C} & T(C)
}\]
This shows $G_T(g\circ_* f) = \mu_C \circ T(\mu_C\circ Tg\circ f) = \mu_C\circ T\mu_C \circ T^2g\circ Tf = (\mu_C\circ Tg)\circ (\mu_B\circ Tf) = G_T g\circ G_T f$.

Conversely, there is a ``coalgebra functor'' $F_T:\mathcal{C}\rightarrow\mathcal{C}_T$, by $F_T(A) = A$, and $F_T f:= \eta_B\circ f:A\rightarrow B\rightarrow T(B)$ for all $A\in\mathcal{C}$ and $f:A\rightarrow B$ in $\mathcal{C}$. This is again functorial, as $F_T \id_A := \eta_A\circ \id_A = \eta_A$ is the identity in $\mathcal{C}_T$, and given $f:A\rightarrow B$, $g:B\rightarrow C$ the following diagram commutes:
\[\xymatrix{
    A \ar[dr]_-{F_T f} \ar[r]^-{f} & B \ar[d]_-{\eta_B} \ar[r]^-{g} & C\ar[d]^-{\eta_C}\\
    & T(B) \ar[dr]_-{T(F_T g)} \ar[r]^-{Tg} & T(C) \ar[d]^-{T\eta_C} \ar[dr]^-{\id_{T(C)}}\\
    && T^2(C) \ar[r]_-{\mu_C} & T(C)
}\]
(Two left triangles are definition of $F_T$, the square is the naturality of $\eta:\Id_C\rightarrow T$, and bottom right is the triangle identity of monad).

This shows $F_T (g\circ f)=\eta_C\circ (g\circ f)=\id_{T(C)}\circ \eta_C\circ g\circ f = \mu_C \circ T(F_T g)\circ F_T f = F_T g\circ_* F_T f$.

Similar to the Eilenberg-Moore category, we have the adjunction:
\begin{thm}
The two functors are adjoint pairs $(F_T\dashv G_T)$, and the induced monad is $(T,\eta,\mu)$.
\end{thm}
\begin{proof}
    Similar manner, we ought to show a natural isomorphism between the two homsets:
    \[\Phi:\Hom_{\mathcal{C}_T}(A,B)\rightleftarrows \Hom_\mathcal{C}(A,T(B)):\Psi\]
    In fact, by definition these two sets are equal, as a morphism $f:A\rightsquigarrow B$ in $\mathcal{C}_T$, is precisely a morphism $f:A\rightarrow T(B)$, so the ``identity map'' between the two sets provide an obvious bijection.

    About functoriality, this is provided by the properties of morphisms that're established beforehand, here we can comfortably omit it :)

    \hfil

    It remains to check they do form the monad $T$: Fix $A\in\mathcal{C}$, and morphism $f:A\rightarrow B$, we get the following description:
    \[G_T F_T (A) = G_T (A) = T(A),\quad G^T(F^T f) = G_T (\eta_B\circ f) = \mu_B \circ T(\eta_B\circ f) = \mu_B\circ T\eta_B\circ Tf = Tf\]
    The morphism equality is due to triangle identity, $\mu_B\circ T\eta_B = \id_{T(B)}$. This shows $T = G_T F_T$.

    About the unit transformation $\eta':\Id_\mathcal{C}\rightarrow G_T F_T$ is given by $\eta'_A = \Phi(\eta_A) =\eta_A$ (because $\eta_A$ is the identity of $A$ in $\mathcal{C}_T$), showing $\eta'=\eta$.

    About the counit transformation $\epsilon':F_T G_T\rightarrow \Id_{\mathcal{C}_T}$, every $B\in\mathcal{C}_T$ satisfies $\epsilon'_B = \Psi(\id_{T(B)}) = \id_{T(B)}:T(B)\rightsquigarrow B$. Then, the multiplication transformation $\mu':=G_T \epsilon_{F_T}$ satisfies $\mu'_A = G_T \epsilon_{F_T (A)} = G_T(\id_{T(A)}) = \mu_A\circ T\id_{T(A)} = \mu_A\circ \id_{T^2(A)} = \mu_A$, showing $\mu' = \mu$.

    So, the given adjunction does induce the monad $(T,\eta,\mu)$, finishing the proof.
\end{proof}

\hfil

This category is also special in the following sense:
\begin{thm}
    The Kleisli Category $\mathcal{C}_T$ with adjoint pair $(F_T\dashv G_T)$ is initial in $\text{Adj}_T$, the category of adjoint pairs inducing $T$.
\end{thm}

Another characterization of Kleisli Category, can be done as a full subcategory of the Eilenberg-Moore category:

\begin{lem}
    The canonical functor $\mathcal{C}_T\rightarrow\mathcal{C}^T$ is fully faithful, with essential image being all the free $T$-algebra (all pairs $(T(A),\mu_A)$).
\end{lem}

Unfortunately for my stamina purpose, I'll omit the proof here :/

This category can also be viewed as the category of all free $T$-algebra. An intuition behind this can be thought of as: Any adjoint pair inducing monad $T$, must contain information of free $T$-algebra (the Kleisli Category), while be contained in the category of all $T$-algebra. Hence, knowing the $T$-algebras also provides all adjunction associated to $T$.


\end{document}